\chapter{Практика}


\section{Настройка интерфейса роутера}

\begin{lstlisting}[language=C, caption={Пример настройки интерфейса роутера}]
#R1
configure

interface gigabitethernet 1/0/1
ip address 192.168.1.1/24
ip firewall disable
exit

do commit
do confirm

\end{lstlisting}


\section{Настройка интерфейса ПК}

\begin{figure}[H]  
    \centering
    \includegraphics[width=0.75\textwidth]{debian.png}
    \caption{Пример настройки Debian}
\end{figure}

\section{Настройка NTP сервера}

Чтобы NTP сервер передавать время на маршрутизатор, нужно, чтобы на самом сервере было настроено время. Для этого можно воспользоваться командами date, чтобы 
посмотреть текущее время и date <MM:DD:HH> для настройки времени. \\

Далее необходимо проверить статус службы NTP с помощью timedatectl status chronyd, а при необходимости включить с помощью timedatectl set-ntp true.\\


\section{Настройка NTP клиента}


\begin{lstlisting}[language=C, caption={Пример настройки NTP на роутере}]
configure

# включаем NTP
ntp enable
# задаем адрес сервера
ntp server 192.168.1.10
# задаем минимальный интервал опроса (2^1 sec)
minpoll 1
# задаем максимальный интервал опроса (2^4 sec)
maxpoll 4

# настройка часового пояса
clock timezone gmt +8

do commit
do confirm


\end{lstlisting}


\section{Проверка работоспособности}

Если всё настроено верно, то время на маршрутизаторе должно синхронизироваться со временем на сервере:


\begin{figure}[H]  
    \centering
    
\includegraphics[width=0.75\textwidth]{time2.png}
    \caption{Время на роутере}
\end{figure}

Время совпадает со временем на сервере, значит, всё настроено верно.

\endinput